% Generated draft: requires mathematical and source review.
\documentclass[11pt]{article}
\usepackage[margin=25mm]{geometry}
\usepackage{fontspec}
\setmainfont{Latin Modern Roman}
\usepackage{amsmath,amssymb,mathtools,unicode-math}
\setmathfont{Latin Modern Math}
\usepackage{enumitem,needspace,xurl,hyperref}
\setlength{\parindent}{0pt}
\setlength{\parskip}{5pt}
\setlength{\emergencystretch}{3em}
\newcommand{\N}{\mathbb{N}}
\newcommand{\Z}{\mathbb{Z}}
\newcommand{\Q}{\mathbb{Q}}
\newcommand{\R}{\mathbb{R}}
\newcommand{\C}{\mathbb{C}}
\newcommand{\F}{\mathbb{F}}
\newcommand{\A}{\mathbb{A}}
\newcommand{\PP}{\mathbb{P}}
\newcommand{\E}{\mathbb{E}}
\newcommand{\eps}{\varepsilon}
\newcommand{\dd}{\,\mathrm{d}}
\begin{document}
\section*{UnsolvedMath ID: 30002046}
\textbf{Source problem code:} catalogue entry
\subsection{Definitions and mathematical statement}
For the standard unit-interval normalization, Minkowski's question-mark function is the map \(\mathord{?}:[0,1]\to[0,1]\) defined by \(\mathord{?}(0)=0\), \(\mathord{?}(1)=1\), and continued fractions as follows. If \(x\in(0,1)\) is irrational, write its regular continued fraction as \(x=[0;a_1,a_2,\ldots]\), where \(a_n\in\Z_{\geq 1}\), and set
\[
\mathord{?}(x)=2\sum_{n=1}^{\infty}(-1)^{n+1}2^{-(a_1+\cdots+a_n)}.
\]
If \(x\in(0,1)\) is rational, use its finite regular continued fraction \(x=[0;a_1,\ldots,a_m]\) with \(a_m\geq2\), and set
\[
\mathord{?}(x)=2\sum_{n=1}^{m}(-1)^{n+1}2^{-(a_1+\cdots+a_n)}.
\]
Under this normalization, the question is to determine the cardinality of the full fixed-point set
\[
\operatorname{Fix}(\mathord{?})=\{x\in[0,1]:\mathord{?}(x)=x\},
\]
including both endpoints and all rational and irrational solutions, without assuming in advance that this set is finite.

The supplied source asks for all real solutions but does not explicitly specify the domain or an extension beyond \([0,1]\). The displayed formulation makes the standard unit-interval interpretation explicit. If a function on a larger real domain is intended, its definition and the corresponding counting domain remain unspecified.
\subsection{Short English statement}
Determine how many fixed points Minkowski's question-mark function has. The precise formulation above uses its standard domain \([0,1]\); the supplied record does not specify an extension to a larger real domain.
\subsection{Sources}
\begin{itemize}
\item Local UnsolvedMath registry, supplied complete record: permanent ID \(30002046\), source problem code \texttt{catalogue entry}, title \emph{Real Fixed Points of the Minkowski Question-Mark Function}, in the Oberwolfach Reports Open Problems collection. The record supplies proposed year 2012 and names no individual proposer.
\item Supplied report citation: \emph{Diophantische Approximationen (2013)}, Oberwolfach Reports, DOI 10.4171/owr/2012/22, pp.\ 1307--1359. \url{https://doi.org/10.4171/owr/2012/22}. The supplied background describes these reports as documenting mathematical workshops and their open-problem discussions. The registry labels the statement \texttt{corrected\_verified} and states that the official report poses this question, with the extraction's corrupted function notation corrected. This source-verification claim is supplied metadata, not an independent verification here.
\item Gayfulin--Shulga, \emph{Diophantine properties of fixed points of Minkowski question mark function}. Supplied literature link: \url{https://arxiv.org/abs/1811.10139}.
\item Supplied reference labelled \emph{Acta Arithmetica version of the fixed-point study}: \url{https://m.mathnet.ru/php/archive.phtml?jrnid=acta&option_lang=rus&paperid=6&wshow=paper}.
\item Supplied literature assessment, checked on exactly \texttt{2026-08-21}: the record labels the problem open and reports that the literature establishes several fixed points and studies their arithmetic, while leaving the total count and the existence of further irrational fixed points undetermined. This is supplied dated curation derived from the OWR material, not an independently verified literature assessment. The supplied background expressly requests independent verification of time-sensitive status and cited claims.
\end{itemize}
\end{document}
